% $Id: inputtop.tex 79414 2026-06-19 15:13:30Z karl $
% Public domain. Originally written by Karl Berry, 2026.
% Test recursive \input behavior. See scan_file_name in tex.ch.
% This is not run as part of make check; it's here for reference.
\input \input inputsub \end % input ourselves

% To run:
%   tex -ini inputrecurse
% Desired output is the error (\relax stops filename scanning):
%   ! I can't find file `'.
% Undesired output is the error (nothing stops scanning, inputsub.tex is read):
%   ! I can't find file `randomtext'.

% Background:
% 2009 - Taco implements \input {...} (group-delimited filenames) for
% LuaTeX, using @<Get the next non-blank non-relax non-call...@>
% and thus ignoring \relax tokens between the \input and {, whether or
% not there is a {.
% commit f17fd9c40931bda4360d54bf7d4d024e79263542
% Author: Taco Hoekwater <taco@elvenkind.com>
% Date:   Tue Nov 3 14:41:38 2009 +0000
% 
% 2020 - Phelype ports the LuaTeX implementation to all other engines,
% via tex.ch.
% 
% 2026 - Udo points out that this allows a recursive \input, contrary
% to Knuth's specification. The outcome is that whatever text begins
% a subsidiary \input file could be scanned as part of a filename.
% Ultimately because
%   @<Initiate or terminate input...@>=
% does
%   else if name_in_progress then insert_relax
% But that inserted \relax is ignored.
% 
% This test file \input s itself instead of a second file, so the
% problematic recursion doesn't stop until TeX's hardwired "input levels"
% is exceeded. 
%
% Taco's original motivation for ignoring \relax was to have the same
% behavior as \toks assignments, which likewise ignore \relax before the
% <balanced text> of the assignment. That's plausible, but there are many
% other instances of TeX syntax where \relax is not ignored before a
% <balanced text>, so this does not seem an overwhelming argument.
%
% While there is no known use for this behavior, a discussion on
% tex-implementors in May--June 2026 could not come to consensus about
% whether the ignoring of \relax should be controlled by a configuration
% variable (just in case) or unconditionally reverted back to the
% traditional behavior, exactly to avoid the complication of creating and
% configuring and documenting and users having to think about yet another
% configuration variable. There are reasonable arguments on both sides.
% 
% As the code in tex.ch stands now, there is no configuration variable.
